%XXXXXXXXXXXXXXXXXXXXXXXXXXXXXXXXXXXXXXXXXXXXXXXXXXXXXXXXXX
%XXXXXXXXXXXXXXXXXXXXXXXXXXXXXXXXXXXXXXXXXXXXXXXXXXXXXXXXXX
%XXXXXXXXXXXXXXXXXXXXXXXXXXXXXXXXXXXXXXXXXXXXXXXXXXXXXXXXXX
\chapter*{Formelzeichen, Abkürzungen und Koordinatensysteme}
%XXXXXXXXXXXXXXXXXXXXXXXXXXXXXXXXXXXXXXXXXXXXXXXXXXXXXXXXXX
%XXXXXXXXXXXXXXXXXXXXXXXXXXXXXXXXXXXXXXXXXXXXXXXXXXXXXXXXXX
%XXXXXXXXXXXXXXXXXXXXXXXXXXXXXXXXXXXXXXXXXXXXXXXXXXXXXXXXXX
\addcontentsline{toc}{chapter}{Formelzeichen, Abkürzungen und
Koordinatensysteme} \markboth{Formelzeichen}{Formelzeichen}

%abkürzender Befehl für das Formelverzeichnis (fvz)
\newcommand\fvz[2]{
  $#1$ & #2 \\
}


\textit{Selten benutzte Formelzeichen und Begriffe sowie abweichende
Bedeutungen werden ausschlie{\ss}lich im Text erläutert. Es werden,
falls nicht gesondert angegeben, SI-Einheiten ($\einh{m}$,
$\einh{kg}$, $\einh{s}$, $\einh{rad}$) verwendet. Verwenden Sie ausschließlich die physikalisch richtige Schreibweise von Einheiten z.\,B. $v/\frac{km}{h}$.}


%%XXXXXXXXXXXXXXXXXXXXXXXXXXXXXXXXXXXXXXXXXXXXXXXXXXXXXXXXXX
\subsection*{Allgemeine Schreibweisen}
%%XXXXXXXXXXXXXXXXXXXXXXXXXXXXXXXXXXXXXXXXXXXXXXXXXXXXXXXXXX
%-----------------------------Tabelle--------------------------------
{\renewcommand{\arraystretch}{1.5} \tabcolsep=0.3cm
\begin{longtable}{>{\RaggedRight}p{20mm} >{\RaggedRight}p{135mm}}
\fvz{a}
 {Variable}

\fvz{\dot{a}}
 {zeitliche Ableitung der Variablen $a$}

\fvz{\vec{a}}
 {Vektor}

\fvz{\V{a}}
 {Spaltenmatrix, Vektor}

\fvz{\M A}
 {Matrix}

\fvz{_K\V{a}_{BC}}
 {Ortsvektor vom Punkt B nach C, angegeben im
 Koordinatensystem K}

\fvz{^{IK}\M{A}}
 {Transformationsmatrix, die Spaltenmatrizen bzw. Vektoren vom
 Koordinatensystem K in das System I transformiert}

\fvz{\hat{a}}
 {geschätzte Grö{\ss}e}

\end{longtable}}
%--------------------------------------------------------------------

\newpage
%XXXXXXXXXXXXXXXXXXXXXXXXXXXXXXXXXXXXXXXXXXXXXXXXXXXXXXXXXX
\subsection*{Lateinische Notation}
%XXXXXXXXXXXXXXXXXXXXXXXXXXXXXXXXXXXXXXXXXXXXXXXXXXXXXXXXXX
%-----------------------------Tabelle--------------------------------
{\renewcommand{\arraystretch}{1.5} \tabcolsep=0.3cm
\begin{longtable}{>{\RaggedRight}p{20mm} >{\RaggedRight}p{135mm}}
\fvz{a}
 {Beschleunigung}

\fvz{b}
 {Spurweite}

\fvz{b_l}
 {Abstand der linken Räder zum Schwerpunkt}

\fvz{b_r}
 {Abstand der rechten Räder zum Schwerpunkt}

\fvz{c}
 {Steifigkeit}

\fvz{h_C}
 {Schwerpunkthöhe über der Fahrbahn}

\fvz{G}
 {übertragungsfunktion}

\fvz{_K\V{F}_G}
 {Vektor der Gewichtskraft im Fahrzeugschwerpunkt, notiert im
 Koordinatensystem $K$}

\fvz{_K\V{F}_{ext}}
 {Vektor der extern angreifenden Kräfte, notiert im
 Koordinatensystem $K$}

\fvz{_RF_{ix}}
 {Kraft am Rad Nr. $i$ in $x$-Richtung des Koordinatensystems $R$, $i = 1,...,4$}

\fvz{_R\V{F}_{i}}
 {Kraftvektor am Rad Nr. $i$ im Koordinatensystem $R$, $i = 1,...,4$}

\fvz{J_{zz}}
 {Massenträgheitsmoment des Fahrzeugs um die Hochachse $z$}

\fvz{l}
 {Radstand}

\fvz{l_h}
 {Abstand der Hinterachse zum Schwerpunkt}

\fvz{l_v}
 {Abstand der Vorderachse zum Schwerpunkt}

\fvz{m}
 {Fahrzeugmasse}

\fvz{_R\V{M}_{i}}
 {Momentenvektor am Rad Nr. $i$ im Koordinatensystem $R$, $i = 1,...,4$}

\fvz{p_i}
 {Bremsdruck am Rad $i$, $i = 1,...,4$}

\fvz{q_i}
 {Kraftschluss bzw. Kraftschlussquotient des Rades Nr. $i$, $i = 1,...,4$}

\fvz{r_{dyn}}
 {dynamischer Reifenradius}

\end{longtable}}
%--------------------------------------------------------------------


%XXXXXXXXXXXXXXXXXXXXXXXXXXXXXXXXXXXXXXXXXXXXXXXXXXXXXXXXXX
\subsection*{Griechische Notation}
%XXXXXXXXXXXXXXXXXXXXXXXXXXXXXXXXXXXXXXXXXXXXXXXXXXXXXXXXXX
%-----------------------------Tabelle--------------------------------
{\renewcommand{\arraystretch}{1.5} \tabcolsep=0.3cm
\begin{longtable}{>{\RaggedRight}p{20mm} >{\RaggedRight}p{135mm}}
\fvz{\alpha}
 {Schräg\-lauf\-winkel}

\fvz{\beta}
 {Schwimmwinkel im Fahrzeugschwerpunkt}

\fvz{\gamma}
 {Sturz}

\fvz{\delta}
 {Lenkwinkel}

\fvz{\V{\epsilon}}
 {Schätzfehler des Zustandsbeobachters}

\fvz{\theta}
 {Nickwinkel}

\fvz{\lambda}
 {Längsschlupf}

\fvz{\mu, \mu_H, \mu_L}
 {Reibwert, Hochreibwert, Niedrigreibwert}

\fvz{\xi_1}
 {Querneigung der Fahrbahn}

\fvz{\xi_2}
 {Längsneigung der Fahrbahn}

\fvz{\tau}
 {Nachlauf}

\fvz{\varphi}
 {Wankwinkel, Ausrichtungswinkel der ebenen Reifenkraft im Bezug
 zur Reifenlängsachse}

\fvz{\psi}
 {Gierwinkel}

\fvz{\omega_i}
 {Kreis- bzw. Winkelgeschwindigkeit am Rad Nr. $i$, $i = 1,...,4$}

\fvz{\zeta_1, \zeta_2}
 {Verkippung des Fahrzeugkörpers gegenüber dem Inertialsystem $I$ in $_Ix$- und $_Iy$-Richtung}

\end{longtable}}
%--------------------------------------------------------------------

\newpage

%XXXXXXXXXXXXXXXXXXXXXXXXXXXXXXXXXXXXXXXXXXXXXXXXXXXXXXXXXX
\subsection*{Abkürzungen}
%XXXXXXXXXXXXXXXXXXXXXXXXXXXXXXXXXXXXXXXXXXXXXXXXXXXXXXXXXX
%-----------------------------Tabelle--------------------------------
{\renewcommand{\arraystretch}{1.5} \tabcolsep=0.3cm
\begin{longtable}{>{\RaggedRight}p{20mm} >{\RaggedRight}p{135mm}}

\fvz{\mbox{ABS}}
  {\textbf{A}nti\textbf{b}lockier\textbf{s}ystem}

\fvz{\mbox{BMW}}
  {\textbf{B}ayrische \textbf{M}otoren \textbf{W}erke}

\fvz{\mbox{C}}
 {Fahrzeugschwerpunkt (engl. "`\textbf{c}enter of gravity"')}

\fvz{\mbox{ESM}}
  {\textbf{E}in\textbf{s}pur\textbf{m}odell}

\fvz{\mbox{FDR}}
  {\textbf{F}ahr\textbf{d}ynamik\textbf{r}egler}

\fvz{\mbox{HR, HL}}
 {\textbf{h}inten \textbf{r}echts, \textbf{h}inten \textbf{l}inks}

\fvz{\mbox{LIDAR}}
 {\textbf{Li}ght \textbf{D}etection \textbf{A}nd \textbf{R}anging,
 Abstandssensor und Umgebungserfassung}

\fvz{\mbox{VR, VL}}
 {\textbf{v}orne \textbf{r}echts, \textbf{v}orne \textbf{l}inks}

\fvz{x}
 {Längsachse des jeweiligen Koordinatensystems}

\fvz{y}
 {Querachse des jeweiligen Koordinatensystems}

\fvz{z}
 {Hochachse des jeweiligen Koordinatensystems}
\end{longtable}}
%--------------------------------------------------------------------




%XXXXXXXXXXXXXXXXXXXXXXXXXXXXXXXXXXXXXXXXXXXXXXXXXXXXXXXXXX
\subsection*{Indizes}
%XXXXXXXXXXXXXXXXXXXXXXXXXXXXXXXXXXXXXXXXXXXXXXXXXXXXXXXXXX
%-----------------------------Tabelle--------------------------------
{\renewcommand{\arraystretch}{1.5} \tabcolsep=0.3cm
\begin{longtable}{>{\RaggedRight}p{20mm} >{\RaggedRight}p{135mm}}
\fvz{0}
 {nomineller Wert/Referenzwert}

\fvz{\mbox{A}}
 {Antrieb}

\fvz{\mbox{B}}
 {Bremse}

\fvz{\mbox{C}}
 {Fahrzeugschwerpunkt (engl. "`center of gravity"')}

\fvz{\mbox{F}}
 {Fahrer}

\fvz{\mbox{HR, HL}}
 {hinten rechts, hinten links}

\fvz{i}
 {Radnummer, $i = 1,2,3,4$ (VR, VL, HL, HR)}

\fvz{j}
 {Achsenindex, $j = v, h$}

\fvz{\mbox{l}}
 {links}

\fvz{\mbox{r}}
 {rechts}

\fvz{\mbox{R}}
 {Regler oder Reifen}

\fvz{\mbox{VR, VL}}
 {vorne rechts, vorne links}

\end{longtable}}
%--------------------------------------------------------------------


%XXXXXXXXXXXXXXXXXXXXXXXXXXXXXXXXXXXXXXXXXXXXXXXXXXXXXXXXXX
\subsection*{Koordinatensysteme}
%XXXXXXXXXXXXXXXXXXXXXXXXXXXXXXXXXXXXXXXXXXXXXXXXXXXXXXXXXX
\textit{Zur Definition der Koordinatensysteme siehe Abbildung
\ref{Bild_TG_Koordinatensysteme} und zur Drehreihenfolge Tabelle
\ref{Tabelle_TG_Koordinatentransformation} auf Seite
\pageref{Tabelle_TG_Koordinatentransformation}.}

%-----------------------------Tabelle--------------------------------
{\renewcommand{\arraystretch}{1.5} \tabcolsep=0.3cm
\begin{longtable}{>{\RaggedRight}p{20mm} >{\RaggedRight}p{135mm}}

\fvz{B}
 {horizontiertes Bahnkoordinatensystem, tangential zur Bahn
 des Fahrzeugschwerpunkts}

\fvz{H}
 {horizontiertes, Fahrzeugkörper festes Koordinatensystem}

\fvz{I}
 {Inertiales, ortsfestes Koordinatensystem}

\fvz{K}
 {körperfestes Fahrzeug-Koordinatensystem}

\fvz{R_i}
 {Radträger festes Radkoordinatensystem des Rads $i$}

\end{longtable}}
%--------------------------------------------------------------------
